\chapter{Le premier mois}\label{ch:premier-mois}

Ce dernier chapitre rassemble, dans l'ordre où on les fait, les gestes qui ont servi sur
tous les fonds. Il vaut pour Leibniz comme pour n'importe quel autre fonds.

\section{Semaine 1 : lire la source}

\begin{itemize}
  \item Remplir la grille du chapitre~\ref{ch:grille} dans une issue, chaque constat daté :
        format, en-têtes, certificat, licence des images, transcriptions existantes,
        édition de référence.
  \item Écrire au détenteur et aux équipes qui travaillent déjà sur le fonds. Garder les
        échanges dans le dépôt privé.
  \item Créer deux dépôts, un public et un privé, chacun avec son \fichier{README} qui dit ce
        qu'il contient et pourquoi. Écrire le \fichier{.gitignore} avant le premier fichier
        (images, PDF, copie de travail, transcriptions en attente).
  \item Écrire le fichier des droits (\fichier{RIGHTS.md}) : ce qui est à soi et sous quelle
        licence, ce qui n'est pas à soi, l'engagement de retrait.
\end{itemize}

\section{Semaine 2 : l'inventaire}

\begin{itemize}
  \item Un script qui lit le catalogue (API, manifestes) et en fait un fichier de données,
        avec la date du relevé. Une pause entre deux requêtes, un \emph{user agent} qui
        nomme le projet, l'arrêt au premier refus.
  \item Une règle de sélection, une seule, écrite (« le volume est servi en entier », « la
        cote est dans LH~35 et n'est pas éditée »).
  \item Les cotes et datations du détenteur, citées comme telles.
\end{itemize}

\section{Semaine 3 : un lot, de bout en bout}

\begin{itemize}
  \item Adapter la procédure écrite de transcription (\fichier{SKILL.md}) à la main et à la
        langue du fonds : abréviations, notation, numérotation des pages.
  \item Transcrire un lot de vingt pages. Le relire soi-même, contre le fac-similé, page
        par page. Noter ce que la procédure n'avait pas prévu, et la corriger.
  \item Afficher le lot sur une page locale, transcription à gauche, fac-similé servi par
        la source à droite.
\end{itemize}

\section{Semaine 4 : ce qu'on publie}

\begin{itemize}
  \item Décider, pour chaque type de fichier, s'il est public, privé ou jamais stocké, et
        l'écrire.
  \item Mettre en ligne le site avec l'inventaire, la méthode, les sources et les droits,
        avant la première transcription si celles-ci attendent un accord.
  \item Vérifier par un script, à chaque déploiement, que la source répond encore comme le
        jour du relevé.
\end{itemize}

\begin{regle}
On ne passe au lot suivant que lorsque le premier a été relu par une personne qui sait lire
la main. Une procédure qui produit vingt lots faux produit vingt fois la même erreur.
\end{regle}

Lean, les trouvailles et les articles viennent ensuite, sur une base dont on sait ce qui a
été vérifié. C'est la seule chose que les cinq fonds de ce livre ont en commun.
